\subsection{Why not propositional logic?}
\label{sec:why-fol}
\slidesrc{FOL: 2--3}

Consider the argument:
\begin{center}
\begin{tabular}{l}
all monitors are ready \\
X12 is a monitor \\
\hline
therefore X12 is ready
\end{tabular}
\end{center}
Intuitively it is \emph{sound}: if we accept that the two premises (above the line) are
true, we must accept that the conclusion is true too. Each of the three statements is a
proposition (true or false), so we could try propositional logic:
\begin{center}
$p$ = ``all monitors are ready'', \quad $q$ = ``X12 is a monitor'', \quad $r$ = ``X12 is ready''.
\end{center}
But then the argument becomes $p,\ q \;/\therefore\; r$, which is \textbf{not} valid: the
valuation $p = q = \true$, $r = \false$ makes the premises true and the conclusion false.
The sense of the argument is lost, because the connection between ``monitors'',
``X12'' and ``ready'' is hidden inside the atomic propositions.

\begin{nbintuition}[Propositional logic is too coarse-grained]
Propositional logic sees sentences as black boxes. The validity of the monitor argument
depends on what is \emph{inside} the boxes: the same individual (X12), the same property
(being a monitor, being ready), and the word ``all''. We need a logic that can name
individuals, express properties and relations, and quantify over all individuals:
first-order logic. In FOL the argument becomes
\[
\forall x\,(\mathit{Monitor}(x) \limp \mathit{Ready}(x)),\quad \mathit{Monitor}(X12)
\;/\therefore\; \mathit{Ready}(X12)
\]
and it is valid (Section~\ref{sec:proof-theory-fol}).
\end{nbintuition}

\begin{nbwhy}
The slides come from a Software Engineering course, where FOL is used for \textbf{formal
specification} (e.g.\ the Z notation): ``every reactor has a temperature between 100 and
1000'' is a requirement about \emph{all} reactors. In AI the same language describes
knowledge bases (``every employee has a manager''), database queries (SQL is essentially
FOL), and the beliefs of agents.
\end{nbwhy}
